%%% Connections: Geometric Foundations
%%% (direction fields, phase portraits, vector fields, flows, manifolds)

\begin{frame}{Geometric Interpretation}
Differential equations define a \textbf{vector field}.

Example:
\[
\frac{dx}{dt} = f(x)
\]

Each point has a vector that determines the direction of motion.

Solutions are called \textbf{integral curves} --- they are trajectories that are everywhere tangent to the field.
\end{frame}


\begin{frame}{Direction Fields}
For a first-order equation
\[
\frac{dy}{dx} = f(x,y)
\]

Each point $(x,y)$ has a slope:
\[
\text{slope} = f(x,y)
\]

Plotting these slopes produces a \textbf{direction field} (or slope field).

Solutions follow these directions --- they are curves that are tangent to the little line segments at every point.
\end{frame}

\begin{frame}{Phase Portraits}
For systems:
\[
\dot{x} = f(x,y), \quad \dot{y} = g(x,y)
\]

We can study trajectories in the $(x,y)$ plane.

This geometric picture is called a \textbf{phase portrait}.

It reveals:
\begin{itemize}
    \item \textbf{Equilibrium points} (where $\dot{x} = \dot{y} = 0$)
    \item \textbf{Stability} (do nearby trajectories approach or flee the equilibrium?)
    \item \textbf{Oscillations} (closed orbits = periodic solutions)
    \item \textbf{Separatrices} (curves that divide qualitatively different behaviours)
\end{itemize}
\end{frame}


\begin{frame}{What Does $\frac{dy}{dx}$ Actually Mean?}
\[
\frac{dy}{dx} = f(x,y)
\]

can be written as a differential form:
\[
dy - f(x,y)\,dx = 0
\]

Solutions are curves along which this 1-form vanishes.

This connects differential equations with
\begin{itemize}
\item differential forms
\item exterior derivatives
\item integrability conditions (Frobenius theorem)
\end{itemize}

\vspace{0.3cm}
\textit{An exact equation $M\,dx + N\,dy = 0$ is one where $\frac{\partial M}{\partial y} = \frac{\partial N}{\partial x}$ --- precisely the condition for the form to be closed.}
\end{frame}

\begin{frame}{Vector Fields}
A first-order equation
\[
\frac{dx}{dt} = f(x)
\]

defines a \textbf{vector field}.

Each point $x$ has an associated velocity vector $f(x)$.

Solutions are curves tangent to the vector field.

These curves are called \textbf{integral curves} --- the same term used in differential geometry for curves whose tangent vectors match a given field.
\end{frame}

\begin{frame}{Flows on Manifolds}
Given a vector field $X$ on a manifold $M$:
\[
\dot{x} = X(x)
\]

its solutions define a \textbf{flow}
\[
\varphi_t : M \rightarrow M
\]

which moves points along integral curves.

Properties of the flow:
\begin{itemize}
    \item $\varphi_0 = \text{id}$ (at time zero, nothing moves)
    \item $\varphi_{t+s} = \varphi_t \circ \varphi_s$ (the group law)
\end{itemize}

Thus differential equations generate \textbf{one-parameter groups of diffeomorphisms} --- dynamical systems on manifolds.
\end{frame}

\begin{frame}{Tangent Vectors}
On a manifold the derivative defines a \textbf{tangent vector}:
\[
\frac{d}{dt}x(t) \in T_xM
\]

where $T_xM$ is the tangent space at $x$.

Differential equations describe curves whose velocities lie in these tangent spaces.

Hence they are naturally formulated using:
\begin{itemize}
\item tangent bundles $TM$
\item vector fields $X : M \to TM$
\item smooth manifolds
\end{itemize}
\end{frame}


\begin{frame}{Vectors in Differential Geometry}
\textbf{Two equivalent definitions} of a tangent vector $v \in T_pM$:

\vspace{0.2cm}
\textbf{1. As an equivalence class of curves:} \\
Two curves $\gamma_1, \gamma_2$ through $p$ are equivalent if they have the same velocity in every chart. Then $v = [\gamma]$ and
\[
v(f) = \frac{d}{dt}f(\gamma(t))\Big|_{t=0}
\]

\textbf{2. As a derivation on functions:} \\
$v : C^\infty(M) \to \mathbb{R}$ is a linear map satisfying the \textbf{Leibniz rule}:
\[
v(fg) = v(f)\,g(p) + f(p)\,v(g)
\]

\vspace{0.2cm}
\textbf{Connection to DEs:} A vector field $X$ assigns a derivation to each point.
An integral curve $\gamma$ of $X$ satisfies $\dot{\gamma}(t) = X(\gamma(t))$ --- i.e., the vector field \emph{is} the ODE, and integral curves \emph{are} the solutions.
\end{frame}


\begin{frame}{One Equation in Many Mathematical Languages}
Consider the differential equation
\[
\frac{dy}{dx} = f(x,y)
\]

The same object can be interpreted in several ways:

\begin{itemize}
\item \textbf{Slope field}: each point $(x,y)$ has slope $f(x,y)$
\item \textbf{Vector field}: $V(x,y) = (1,\, f(x,y))$
\item \textbf{Differential form}: $dy - f(x,y)\,dx = 0$
\item \textbf{Integral curves}: solutions are curves tangent to $V$
\item \textbf{Flow of a vector field}: trajectories generated by motion in the field
\item \textbf{Dynamical system}: evolution rule for a state
\end{itemize}

A differential equation is therefore a \textbf{geometric object}, not merely an algebraic one.
\end{frame}


\begin{frame}{Differential Equations as Geometry}

A hierarchy of ideas:
\[
\text{functions} \xrightarrow{\;\text{differentiate}\;} \text{derivatives} \xrightarrow{\;\text{assign to points}\;} \text{vector fields}
\]
\[
\xrightarrow{\;\text{integrate}\;} \text{flows} \xrightarrow{\;\text{study globally}\;} \text{dynamical systems}
\]

\vspace{0.2cm}

These systems live on geometric \textbf{spaces}:
\begin{itemize}
\item $\mathbb{R}^n$ (the familiar setting)
\item Smooth manifolds (curved spaces)
\item Configuration spaces in mechanics (position + angle)
\item Phase spaces (position + momentum)
\end{itemize}

\vspace{0.3cm}
This viewpoint connects differential equations with \textbf{geometry, topology, and physics}.
\end{frame}


\begin{frame}{From Curves to Manifolds}
A curve $x(t)$ has velocity $\dot{x}(t)$,
which is a vector in the \textbf{tangent space}:
\[
\dot{x}(t) \in T_{x(t)}M
\]

Thus a differential equation $\dot{x} = X(x)$ specifies a \textbf{section of the tangent bundle}:
\[
X : M \rightarrow TM
\]

and solutions are curves whose velocity matches the field at every point.

This is the geometric formulation of differential equations on manifolds --- it works in curved spaces where coordinates are only local.
\end{frame}

